%==============================================================================
%  ATB ATM LOCATOR -- Technical Research Paper / Article de recherche technique
%  Bilingual edition : Version Francaise + English Version
%  Compile with : pdflatex (run twice for the table of contents / references)
%  Tested with a standard TeX Live / MiKTeX distribution.
%==============================================================================
\documentclass[11pt,a4paper]{article}

%------------------------------------------------------------------ Encoding
\usepackage[T1]{fontenc}
\usepackage[utf8]{inputenc}

%------------------------------------------------------------------ Languages
% The last declared language is the default; we switch explicitly per \part.
\usepackage[english,french]{babel}

%------------------------------------------------------------------ Packages
\usepackage{geometry}
\geometry{margin=2.5cm}
\usepackage{graphicx}
\usepackage{booktabs}
\usepackage{array}
\usepackage{enumitem}
\usepackage{xcolor}
\usepackage{amssymb}
\usepackage{hyperref}
\hypersetup{
  colorlinks=true,
  linkcolor=[HTML]{8B1832},
  urlcolor=[HTML]{8B1832},
  citecolor=[HTML]{8B1832},
  pdftitle={ATB ATM Locator -- Technical Choices and Architecture},
  pdfauthor={Hamza Achour}
}

%------------------------------------------------------------------ Helpers
\definecolor{atbcrimson}{HTML}{8B1832}
\newcommand{\code}[1]{\texttt{#1}}

%==============================================================================
\title{\textbf{Choix Techniques et Architecture d'une Application Mobile\\
de Geolocalisation Bancaire}\\[4pt]
\large Etude de cas du projet \textit{ATB ATM Locator}\\[10pt]
\textbf{Technical Choices and Architecture of a Geolocation-Based\\
Mobile Banking Application}\\[4pt]
\large A Case Study of the \textit{ATB ATM Locator} Project}

\author{
  Hamza Achour\\
  \small Senior Mobile Software Engineer \& Independent Researcher\\
  \small Arab Tunisian Bank --- Mobile Innovation (Case Study)\\
  \small \texttt{1hamza.achour@gmail.com}
}
\date{\today}

%==============================================================================
\begin{document}
\maketitle
\thispagestyle{empty}

\begin{center}
\fbox{\begin{minipage}{0.86\textwidth}\centering\small
\textit{Document bilingue --- la version fran\c{c}aise est pr\'esent\'ee en premier
(\textbf{Partie I}), suivie de la traduction anglaise int\'egrale (\textbf{Partie II}).}\\[2pt]
\textit{Bilingual document --- the French version is presented first
(\textbf{Part~I}), immediately followed by the complete English translation
(\textbf{Part~II}).}
\end{minipage}}
\end{center}

\vspace{6pt}
\tableofcontents
\newpage

%##############################################################################
%  PART I -- FRENCH VERSION
%##############################################################################
\part{Version Fran\c{c}aise}
\selectlanguage{french}

\begin{abstract}
\noindent
\textbf{ATB ATM Locator} est une application mobile bancaire con\c{c}ue pour
permettre aux clients de l'Arab Tunisian Bank de localiser en temps r\'eel les
distributeurs automatiques de billets (DAB) et les agences les plus proches,
d'obtenir un itin\'eraire fiable et de consulter ces informations m\^eme en
l'absence de connexion r\'eseau. Cet article pr\'esente, du point de vue d'un
ing\'enieur logiciel mobile senior, les \emph{choix techniques} structurants du
projet ainsi que leur justification. Les principaux d\'efis d'ing\'enierie
trait\'es sont~: la r\'ecup\'eration et le tri par proximit\'e de donn\'ees
g\'eospatiales h\'et\'erog\`enes, le rendu cartographique performant, une
strat\'egie \emph{offline-first} fond\'ee sur la mise en cache de tuiles, la
gestion fine de la consommation de batterie li\'ee au GPS, et la s\'ecurisation
de la couche r\'eseau. La pile technologique retenue repose sur
\textbf{Flutter/Dart} pour l'interface multiplateforme, \textbf{flutter\_map}
adoss\'e \`a \textbf{OpenStreetMap} pour la cartographie, l'\textbf{API Overpass}
comme source de donn\'ees DAB, le paquet \textbf{geolocator} pour la
g\'eolocalisation, et un cache hors-ligne au format \textbf{MBTiles}
(\textbf{SQLite}). Nous montrons qu'une architecture en couches, coupl\'ee \`a
des m\'ecanismes syst\'ematiques de d\'egradation gracieuse, permet d'atteindre
une robustesse de niveau production tout en restant simple \`a maintenir et
\'economiquement viable (absence de cl\'es propri\'etaires payantes).
\end{abstract}

\section{Introduction}
\subsection{Contexte du projet}
La densit\'e du r\'eseau de DAB en Tunisie et l'h\'et\'erog\'en\'eit\'e des
horaires de service rendent fr\'equente une question simple mais r\'ecurrente
pour l'usager~: \emph{o\`u se trouve le distributeur ouvert le plus proche, et
propose-t-il le service de d\'ep\^ot ou de retrait sans carte~?} \textbf{ATB ATM
Locator} r\'epond \`a ce besoin en combinant une carte interactive, un tri par
distance g\'eod\'esique et un syst\`eme de filtres compos\'es.

\subsection{Probl\'ematique}
La construction d'un outil bancaire fond\'e sur la g\'eolocalisation soul\`eve
plusieurs probl\`emes d'ing\'enierie non triviaux~:
\begin{itemize}[leftmargin=1.4em]
  \item \textbf{Suivi en temps r\'eel}~: agr\'eger des donn\'ees de DAB
  potentiellement incompl\`etes et les ordonner selon la position courante de
  l'utilisateur.
  \item \textbf{Disponibilit\'e hors-ligne}~: garantir l'usage de la carte dans
  des zones \`a faible couverture r\'eseau.
  \item \textbf{Routage fiable}~: d\'el\'eguer la navigation \`a un fournisseur
  d'itin\'eraires \'eprouv\'e sans capturer de donn\'ees sensibles.
  \item \textbf{Sobri\'et\'e \'energ\'etique}~: limiter l'impact du GPS sur
  l'autonomie de l'appareil.
  \item \textbf{S\'ecurit\'e}~: prot\'eger les secrets applicatifs et minimiser
  la surface d'attaque en l'absence de backend propri\'etaire.
\end{itemize}

\subsection{Objectifs de l'architecture technique}
L'architecture vise quatre propri\'et\'es cardinales~: la
\textbf{portabilit\'e} (un code source unique pour Android et iOS), la
\textbf{r\'esilience} (aucune fonctionnalit\'e critique ne d\'epend d'un unique
point de d\'efaillance), la \textbf{testabilit\'e} (logique m\'etier isol\'ee de
l'interface) et la \textbf{soutenabilit\'e \'economique} (recours exclusif \`a
des services gratuits ou \`a offre gratuite).

\section{\'Etat de l'art}
\subsection{Paradigmes de d\'eveloppement mobile}
Trois grandes familles de paradigmes coexistent pour une application fortement
cartographique et orient\'ee localisation~: le d\'eveloppement \textbf{natif}
(Kotlin/Swift), les approches \textbf{multiplateformes \`a moteur de rendu
propre} (Flutter), et les approches \textbf{multiplateformes \`a pont JavaScript}
(React Native). Le natif offre l'acc\`es le plus direct aux capteurs et la
meilleure int\'egration syst\`eme, au prix d'un double effort de
d\'eveloppement. React Native mutualise la logique mais d\'epend d'un pont vers
les composants natifs, ce qui peut p\'enaliser les animations cartographiques
intensives. Flutter compile en code natif ARM et dessine son interface via le
moteur \textbf{Skia/Impeller}, offrant un contr\^ole pixel-pr\`es du rendu,
particuli\`erement avantageux pour des surcouches de marqueurs personnalis\'es.

\subsection{Solutions cartographiques}
Le march\'e propose principalement le \textbf{SDK Google Maps}, \textbf{Mapbox},
et l'\'ecosyst\`eme \textbf{OpenStreetMap} (OSM) rendu via des biblioth\`eques
de type \emph{Leaflet} (dont \code{flutter\_map} est l'analogue Flutter). Le
crit\`ere discriminant pour ce projet est le co\^ut marginal nul et l'absence de
cl\'e propri\'etaire d'OSM, conjugu\'es \`a la disponibilit\'e d'une source de
donnees DAB ouverte (Overpass) align\'ee sur le m\^eme r\'ef\'erentiel
g\'eographique. Le Tableau~\ref{tab:maps-fr} synth\'etise cette comparaison.

\section{Choix Techniques et Architecture}
\subsection{S\'election du framework et du langage}
\textbf{Flutter 3.44} (langage \textbf{Dart 3.12}) a \'et\'e retenu. Les
justifications principales sont~:
\begin{itemize}[leftmargin=1.4em]
  \item \textbf{Code source unique} compil\'e pour Android et iOS, divisant par
  deux l'effort de maintenance.
  \item \textbf{Rendu Skia} ind\'ependant des widgets natifs, garantissant un
  comportement identique des couches de marqueurs sur toutes les plateformes.
  \item \textbf{Hot reload} acc\'el\'erant l'it\'eration sur l'interface
  cartographique.
  \item \textbf{\'Ecosyst\`eme de paquets} mature pour la g\'eolocalisation, la
  base SQLite et le t\'el\'echargement HTTP.
\end{itemize}
Le Tableau~\ref{tab:framework-fr} compare les options \'evalu\'ees.

\begin{table}[ht]
\centering
\caption{Comparaison des paradigmes de d\'eveloppement \'evalu\'es.}
\label{tab:framework-fr}
\small
\begin{tabular}{>{\raggedright\arraybackslash}p{2.6cm} >{\raggedright\arraybackslash}p{5.4cm} >{\raggedright\arraybackslash}p{5.4cm}}
\toprule
\textbf{Technologie} & \textbf{Avantages} & \textbf{Inconv\'enients} \\
\midrule
\textbf{Flutter\,/\,Dart} \textit{(retenu)} &
Code unique multiplateforme~; rendu Skia constant~; excellentes performances de
surcouche cartographique~; hot reload &
Taille du binaire plus \'elev\'ee~; \'ecosyst\`eme bancaire natif moins fourni \\
\midrule
\textbf{React Native} &
Large communaut\'e~; partage de logique JS~; nombreux modules &
Pont JS p\'enalisant le rendu intensif~; fragmentation des biblioth\`eques de
cartes \\
\midrule
\textbf{Natif (Kotlin / Swift)} &
Acc\`es capteur optimal~; int\'egration syst\`eme maximale &
Double base de code~; co\^ut et d\'elai de d\'eveloppement doubl\'es \\
\bottomrule
\end{tabular}
\end{table}

\subsection{Architecture de gestion de l'\'etat}
Le projet adopte une approche d\'elib\'er\'ement \textbf{l\'eg\`ere}~: un \'etat
local r\'eactif (\code{StatefulWidget} + \code{setState}) au niveau des \'ecrans,
adoss\'e \`a une \textbf{couche de services sans \'etat} (classes \`a m\'ethodes
statiques~: \code{ATMService}, \code{ChatService}, \code{LocationService},
\code{OfflineMapService}). Ce choix se justifie par la \textbf{cardinalit\'e
mod\'er\'ee de l'\'etat partag\'e}~: chaque \'ecran poss\`ede un \'etat
relativement autonome, et la navigation par onglets repose sur un
\code{IndexedStack} qui pr\'eserve l'\'etat sans reconstruction.

\subsubsection{Alternatives consid\'er\'ees}
Les solutions \code{Provider}, \code{Riverpod} et \code{BLoC} ont \'et\'e
\'evalu\'ees. \code{BLoC} impose une c\'er\'emonie (\'ev\'enements/\'etats)
disproportionn\'ee pour la taille actuelle du domaine. \code{Riverpod}
constitue la \textbf{trajectoire de migration recommand\'ee} d\`es lors que
l'application int\'egrera un v\'eritable backend transactionnel et un \'etat
d'authentification global. \`A ce stade, \code{setState} maximise la lisibilit\'e
et minimise la dette conceptuelle.

\subsection{G\'eolocalisation, SDK cartographiques et strat\'egie offline-first}
\subsubsection{G\'eolocalisation}
Le paquet \code{geolocator} encapsule la demande de permission, la v\'erification
de l'activation du service et l'acquisition de la position avec
\code{LocationAccuracy.high}. La distance utilisateur--DAB est calcul\'ee via la
formule g\'eod\'esique fournie par \code{Geolocator.distanceBetween}, puis sert
de clef de tri croissant.

\subsubsection{Cartographie}
Le rendu s'appuie sur \code{flutter\_map} (\'equivalent Flutter de Leaflet) et
les tuiles raster d'\textbf{OpenStreetMap}. Les marqueurs sont color\'es selon
l'\'etat du DAB~: bordeaux (ouvert), gris (ferm\'e), contour ambre
(s\'electionn\'e). La navigation routi\`ere est \textbf{d\'el\'egu\'ee} \`a
Google Maps par un lien profond (\code{url\_launcher}), \'evitant ainsi toute
d\'ependance \`a un SDK de routage propri\'etaire.

\subsubsection{Donn\'ees DAB et d\'egradation gracieuse}
Les distributeurs sont r\'ecup\'er\'es aupr\`es de l'\textbf{API Overpass} par
une requ\^ete born\'ee (\code{amenity=atm}, rayon param\'etrable, d\'elai
d'expiration de 15~s). En cas d'\'echec (refus GPS, latence, indisponibilit\'e),
l'application bascule sur un \textbf{jeu de 20 DAB tunisiens pr\'echarg\'es}~: la
carte n'est \emph{jamais} vide.

\subsubsection{Strat\'egie offline-first}
Une r\'egion (p.~ex. \textit{Grand Tunis}) peut \^etre t\'el\'echarg\'ee sous
forme de fichier \textbf{MBTiles}, c'est-\`a-dire une base \textbf{SQLite}
contenant les tuiles PNG. Les coordonn\'ees de tuiles sont calcul\'ees par la
projection \emph{slippy map} et stock\'ees en ordre \textbf{TMS}. Un fournisseur
sur-mesure, \code{MBTilesTileProvider}, fait le pont entre SQLite et
\code{flutter\_map}, rendant la carte exploitable \textbf{sans aucune
connexion}.

\begin{table}[ht]
\centering
\caption{Comparaison des solutions cartographiques.}
\label{tab:maps-fr}
\small
\begin{tabular}{>{\raggedright\arraybackslash}p{3.0cm} >{\raggedright\arraybackslash}p{5.2cm} >{\raggedright\arraybackslash}p{5.2cm}}
\toprule
\textbf{Solution} & \textbf{Avantages} & \textbf{Inconv\'enients} \\
\midrule
\textbf{flutter\_map + OSM} \textit{(retenu)} &
Gratuit, sans cl\'e API~; contr\^ole total du rendu~; coh\'erence avec les
donn\'ees Overpass~; cache MBTiles ais\'e &
Quotas d'usage \'equitable OSM~; pas de trafic temps r\'eel natif \\
\midrule
\textbf{SDK Google Maps} &
Donn\'ees riches~; trafic et POI~; rendu optimis\'e &
Cl\'e payante au-del\`a d'un seuil~; conditions restrictives~; cache hors-ligne
limit\'e \\
\midrule
\textbf{Mapbox} &
Rendu vectoriel \'el\'egant~; offline officiel &
Mod\`ele tarifaire \`a l'usage~; d\'ependance \`a un fournisseur \\
\bottomrule
\end{tabular}
\end{table}

\subsection{Couche r\'eseau, communication API et s\'ecurit\'e}
La couche r\'eseau combine \code{http} (appels REST simples vers Overpass et
Groq) et \code{Dio} (t\'el\'echargement de tuiles en flux d'octets avec
\code{CancelToken}). Toutes les communications transitent en \textbf{HTTPS/TLS}.
Les requ\^etes sont prot\'eg\'ees par des d\'elais d'expiration explicites et
encapsul\'ees dans des blocs de capture d'erreur qui d\'eclenchent la
d\'egradation gracieuse.

\subsubsection{Gestion des secrets}
La cl\'e d'API est isol\'ee dans un fichier \code{secrets.dart} \textbf{exclu du
suivi de version} (\code{.gitignore}), accompagn\'e d'un mod\`ele
\code{secrets.example.dart} versionn\'e. Pour un d\'eploiement de production,
nous recommandons de \textbf{d\'eporter l'appel sensible derri\`ere un proxy
backend} afin de ne jamais embarquer de secret dans le binaire distribu\'e, et
d'injecter la configuration via \code{--dart-define} au moment de la
compilation.

\subsubsection{Principe de moindre privil\`ege}
L'application ne demande que les permissions \code{INTERNET} et de localisation,
ne stocke aucune donn\'ee personnelle identifiable, et conserve la base de
tuiles dans le r\'epertoire priv\'e de l'application, supprimant le besoin de
permissions de stockage externe.

\begin{figure}[ht]
\centering
\fbox{\begin{minipage}{0.86\textwidth}\centering\vspace{6pt}
\textbf{[ DIAGRAMME D'ARCHITECTURE --- ESPACE R\'ESERV\'E ]}\\[10pt]
\fbox{Couche Pr\'esentation : \'Ecrans \& Widgets (Accueil, Carte, Assistance\dots)}\\[4pt]
$\downarrow$\\[4pt]
\fbox{Couche \'Etat : \code{setState} + \code{IndexedStack}}\\[4pt]
$\downarrow$\\[4pt]
\fbox{Couche Services : ATM \textbullet{} Chat \textbullet{} Location \textbullet{} OfflineMap}\\[4pt]
$\downarrow$\\[4pt]
\fbox{Sources de donn\'ees : API Overpass \textbullet{} API Groq \textbullet{} MBTiles/SQLite \textbullet{} GPS}
\vspace{6pt}
\end{minipage}}
\caption{Architecture en couches de \textit{ATB ATM Locator} (sch\'ema de principe).}
\label{fig:arch-fr}
\end{figure}

\section{D\'efis d'impl\'ementation}
\subsection{Consommation de batterie li\'ee au GPS}
L'acquisition continue de position \'epuise rapidement la batterie. La strat\'egie
retenue privil\'egie une \textbf{acquisition ponctuelle} (\code{getCurrentPosition})
plut\^ot qu'un flux permanent, l'\'etat statique des DAB ne justifiant pas un
suivi temps r\'eel. Ce compromis r\'eduit l'empreinte \'energ\'etique tout en
conservant une pr\'ecision \'elev\'ee au moment du calcul des distances.

\subsection{Pr\'ecision et lisibilit\'e des marqueurs}
La superposition de nombreux marqueurs sur une zone dense pose un probl\`eme de
lisibilit\'e. La conception actuelle distingue visuellement la s\'election
(agrandissement et contour ambre) et conserve un nombre de marqueurs ma\^itris\'e
gr\^ace au filtrage compos\'e. L'\textbf{agr\'egation (clustering)} de marqueurs
constitue l'\'evolution naturelle pour les densit\'es \'elev\'ees.

\subsection{Conformit\'e \`a la politique d'usage d'OSM}
Le t\'el\'echargement massif de tuiles doit respecter la politique d'usage
\'equitable d'OSM. L'impl\'ementation applique une \textbf{temporisation de
50\,ms} entre tuiles et un \textbf{User-Agent} identifiant explicitement
l'application, pr\'evenant tout blocage c\^ot\'e serveur.

\subsection{T\'el\'echargements reprenables}
Une coupure r\'eseau ne doit pas invalider un t\'el\'echargement partiel. Chaque
tuile d\'ej\`a pr\'esente en base est d\'etect\'ee et ignor\'ee~: le
t\'el\'echargement reprend exactement o\`u il s'\'etait interrompu, gr\^ace \`a
un point de contr\^ole au niveau de la tuile (\emph{checkpointing}).

\subsection{Conversion de coordonn\'ees et cycle de vie asynchrone}
La conversion entre l'ordonn\'ee XYZ et l'ordonn\'ee \textbf{TMS} (inversion de
l'axe~$y$) est centralis\'ee pour \'eviter les erreurs d'alignement. Par
ailleurs, les rappels asynchrones v\'erifient syst\'ematiquement l'\'etat
\code{mounted} du widget, et un drapeau de d\'emarrage unique \'evite le
re-d\'eclenchement d'un t\'el\'echargement \`a chaque reconstruction d'un
\code{StatefulBuilder}.

\section{Performances et CI/CD}
\subsection{Optimisation du rendu et gestion m\'emoire}
Plusieurs leviers sont mobilis\'es~: l'usage d'\code{IndexedStack} \'evite la
reconstruction des onglets~; l'emploi de constructeurs \code{const} r\'eduit les
reconstructions superflues~; la base de tuiles est ouverte en \textbf{lecture
seule} lors de la consultation et \textbf{ferm\'ee} explicitement dans
\code{dispose}~; les contr\^oleurs (carte, d\'efilement, animation) sont
lib\'er\'es pour pr\'evenir les fuites m\'emoire. Le stockage des tuiles sous
forme de BLOB dans SQLite \'evite la fragmentation du syst\`eme de fichiers.

\subsection{Cha\^ine d'int\'egration et de d\'eploiement continu}
La cha\^ine de livraison recommand\'ee s'articule autour de \textbf{GitHub
Actions}~:
\begin{itemize}[leftmargin=1.4em]
  \item \textbf{Analyse statique}~: \code{flutter analyze} avec
  \code{flutter\_lints} \`a chaque \emph{pull request}.
  \item \textbf{Tests}~: ex\'ecution de \code{flutter test} (widgets et logique
  des services).
  \item \textbf{Construction}~: \code{flutter build apk\,/\,appbundle} et
  \code{ipa}, secrets inject\'es via \code{--dart-define}.
  \item \textbf{Signature et publication}~: automatisation via \textbf{fastlane}
  vers les canaux de test, puis d\'eploiement progressif (\emph{staged
  rollout}).
\end{itemize}

\section{Conclusion}
L'architecture de \textbf{ATB ATM Locator} d\'emontre qu'une s\'eparation nette
en couches, combin\'ee \`a une d\'egradation gracieuse syst\'ematique, suffit \`a
atteindre une robustesse de niveau production sans recourir \`a des services
propri\'etaires co\^uteux. La \textbf{soutenabilit\'e} repose sur l'usage
exclusif de briques ouvertes (OSM, Overpass) et la \textbf{p\'erennit\'e} sur la
modularit\'e de la couche de services, qui permettra de substituer un backend
transactionnel, d'introduire \code{Riverpod} pour un \'etat global,
d'ajouter l'authentification biom\'etrique et les notifications \emph{push}.
La trajectoire d'\'evolution est ainsi clairement balis\'ee, ce qui qualifie
l'architecture de \emph{future-proof} \`a l'\'echelle d'un produit bancaire
grand public.

%##############################################################################
%  PART II -- ENGLISH VERSION
%##############################################################################
\newpage
\part{English Version}
\selectlanguage{english}

\begin{abstract}
\noindent
\textbf{ATB ATM Locator} is a mobile banking application designed to let Arab
Tunisian Bank customers locate, in real time, the nearest Automated Teller
Machines (ATMs) and branches, obtain reliable routing, and consult this
information even without network connectivity. From the standpoint of a senior
mobile software engineer, this paper presents and justifies the project's
defining \emph{technical choices}. The core engineering challenges addressed
are: retrieving and proximity-sorting heterogeneous geospatial data, performant
map rendering, an \emph{offline-first} strategy based on tile caching,
fine-grained control of GPS-induced battery drain, and securing the network
layer. The selected stack relies on \textbf{Flutter/Dart} for the
cross-platform interface, \textbf{flutter\_map} backed by \textbf{OpenStreetMap}
for mapping, the \textbf{Overpass API} as the ATM data source, the
\textbf{geolocator} package for location, and an offline cache in the
\textbf{MBTiles} (\textbf{SQLite}) format. We show that a layered architecture,
coupled with systematic graceful-degradation mechanisms, achieves
production-grade robustness while remaining simple to maintain and economically
sustainable (no paid proprietary keys).
\end{abstract}

\section{Introduction}
\subsection{Project context}
The density of the Tunisian ATM network and the heterogeneity of service hours
make a simple but recurring question common for users: \emph{where is the
nearest open ATM, and does it offer the deposit or cardless-withdrawal service?}
\textbf{ATB ATM Locator} answers this need by combining an interactive map,
geodesic distance sorting, and a composable filter system.

\subsection{Problem statement}
Building a geolocation-based banking tool raises several non-trivial engineering
problems:
\begin{itemize}[leftmargin=1.4em]
  \item \textbf{Real-time tracking}: aggregating potentially incomplete ATM data
  and ordering it by the user's current position.
  \item \textbf{Offline availability}: guaranteeing map usage in areas with poor
  network coverage.
  \item \textbf{Reliable routing}: delegating navigation to a proven routing
  provider without capturing sensitive data.
  \item \textbf{Energy frugality}: limiting the impact of GPS on device battery
  life.
  \item \textbf{Security}: protecting application secrets and minimizing the
  attack surface in the absence of a proprietary backend.
\end{itemize}

\subsection{Objectives of the technical architecture}
The architecture targets four cardinal properties: \textbf{portability} (a single
source base for Android and iOS), \textbf{resilience} (no critical feature
depends on a single point of failure), \textbf{testability} (business logic
isolated from the interface), and \textbf{economic sustainability} (exclusive
reliance on free or free-tier services).

\section{State of the Art}
\subsection{Mobile development paradigms}
Three broad paradigm families coexist for a map-heavy, location-oriented
application: \textbf{native} development (Kotlin/Swift), \textbf{cross-platform
with a dedicated rendering engine} (Flutter), and \textbf{cross-platform with a
JavaScript bridge} (React Native). Native offers the most direct sensor access
and the best system integration, at the cost of duplicated development effort.
React Native shares logic but depends on a bridge to native components, which can
penalize intensive map animations. Flutter compiles to native ARM code and
paints its interface through the \textbf{Skia/Impeller} engine, providing
pixel-accurate rendering control that is particularly advantageous for custom
marker overlays.

\subsection{Mapping solutions}
The market mainly offers the \textbf{Google Maps SDK}, \textbf{Mapbox}, and the
\textbf{OpenStreetMap} (OSM) ecosystem rendered through \emph{Leaflet}-style
libraries (of which \code{flutter\_map} is the Flutter analogue). The
discriminating criterion for this project is OSM's zero marginal cost and the
absence of a proprietary key, combined with the availability of an open ATM data
source (Overpass) aligned with the same geographic reference. Table~%
\ref{tab:maps-en} summarizes this comparison.

\section{Technical Choices and Architecture}
\subsection{Framework and language selection}
\textbf{Flutter 3.44} (\textbf{Dart 3.12} language) was selected. The main
justifications are:
\begin{itemize}[leftmargin=1.4em]
  \item \textbf{Single source base} compiled for Android and iOS, halving
  maintenance effort.
  \item \textbf{Skia rendering} independent of native widgets, guaranteeing
  identical behaviour of marker layers across platforms.
  \item \textbf{Hot reload} accelerating iteration on the map interface.
  \item \textbf{Mature package ecosystem} for geolocation, the SQLite database,
  and HTTP downloading.
\end{itemize}
Table~\ref{tab:framework-en} compares the evaluated options.

\begin{table}[ht]
\centering
\caption{Comparison of the evaluated development paradigms.}
\label{tab:framework-en}
\small
\begin{tabular}{>{\raggedright\arraybackslash}p{2.6cm} >{\raggedright\arraybackslash}p{5.4cm} >{\raggedright\arraybackslash}p{5.4cm}}
\toprule
\textbf{Technology} & \textbf{Pros} & \textbf{Cons} \\
\midrule
\textbf{Flutter\,/\,Dart} \textit{(selected)} &
Single cross-platform codebase; consistent Skia rendering; excellent map-overlay
performance; hot reload &
Larger binary size; less mature native banking ecosystem \\
\midrule
\textbf{React Native} &
Large community; shared JS logic; many modules &
JS bridge penalizes intensive rendering; fragmented map libraries \\
\midrule
\textbf{Native (Kotlin / Swift)} &
Optimal sensor access; maximal system integration &
Dual codebase; doubled development cost and time \\
\bottomrule
\end{tabular}
\end{table}

\subsection{State management architecture}
The project adopts a deliberately \textbf{lightweight} approach: reactive local
state (\code{StatefulWidget} + \code{setState}) at the screen level, backed by a
\textbf{stateless service layer} (classes with static methods:
\code{ATMService}, \code{ChatService}, \code{LocationService},
\code{OfflineMapService}). This choice is justified by the \textbf{moderate
cardinality of shared state}: each screen owns a relatively autonomous state, and
tab navigation relies on an \code{IndexedStack} that preserves state without
rebuilding.

\subsubsection{Considered alternatives}
The \code{Provider}, \code{Riverpod}, and \code{BLoC} solutions were evaluated.
\code{BLoC} imposes an event/state ceremony disproportionate to the current
domain size. \code{Riverpod} is the \textbf{recommended migration path} once the
application integrates a genuine transactional backend and a global
authentication state. At this stage, \code{setState} maximizes readability and
minimizes conceptual debt.

\subsection{Geolocation, map SDKs, and offline-first strategy}
\subsubsection{Geolocation}
The \code{geolocator} package encapsulates permission requests, service-enabled
checks, and position acquisition with \code{LocationAccuracy.high}. The
user--ATM distance is computed via the geodesic helper
\code{Geolocator.distanceBetween} and then serves as the ascending sort key.

\subsubsection{Mapping}
Rendering relies on \code{flutter\_map} (the Flutter equivalent of Leaflet) and
\textbf{OpenStreetMap} raster tiles. Markers are colour-coded by ATM state:
crimson (open), grey (closed), amber outline (selected). Road navigation is
\textbf{delegated} to Google Maps through a deep link (\code{url\_launcher}),
thereby avoiding any dependency on a proprietary routing SDK.

\subsubsection{ATM data and graceful degradation}
ATMs are fetched from the \textbf{Overpass API} via a bounded query
(\code{amenity=atm}, configurable radius, 15\,s timeout). On failure (GPS
denial, latency, unavailability), the application falls back to a
\textbf{preloaded set of 20 Tunisian ATMs}: the map is \emph{never} empty.

\subsubsection{Offline-first strategy}
A region (e.g.\ \textit{Greater Tunis}) can be downloaded as an \textbf{MBTiles}
file, i.e.\ a \textbf{SQLite} database containing PNG tiles. Tile coordinates are
computed by the \emph{slippy map} projection and stored in \textbf{TMS} order. A
custom provider, \code{MBTilesTileProvider}, bridges SQLite and
\code{flutter\_map}, making the map usable \textbf{without any connection}.

\begin{table}[ht]
\centering
\caption{Comparison of mapping solutions.}
\label{tab:maps-en}
\small
\begin{tabular}{>{\raggedright\arraybackslash}p{3.0cm} >{\raggedright\arraybackslash}p{5.2cm} >{\raggedright\arraybackslash}p{5.2cm}}
\toprule
\textbf{Solution} & \textbf{Pros} & \textbf{Cons} \\
\midrule
\textbf{flutter\_map + OSM} \textit{(selected)} &
Free, no API key; full rendering control; consistency with Overpass data; easy
MBTiles caching &
OSM fair-use quotas; no native real-time traffic \\
\midrule
\textbf{Google Maps SDK} &
Rich data; traffic and POIs; optimized rendering &
Paid key beyond a threshold; restrictive terms; limited offline caching \\
\midrule
\textbf{Mapbox} &
Elegant vector rendering; official offline support &
Usage-based pricing; vendor lock-in \\
\bottomrule
\end{tabular}
\end{table}

\subsection{Network layer, API communication, and security}
The network layer combines \code{http} (simple REST calls to Overpass and Groq)
and \code{Dio} (byte-stream tile downloading with \code{CancelToken}). All
communication travels over \textbf{HTTPS/TLS}. Requests are protected by
explicit timeouts and wrapped in error-capturing blocks that trigger graceful
degradation.

\subsubsection{Secrets management}
The API key is isolated in a \code{secrets.dart} file \textbf{excluded from
version control} (\code{.gitignore}), accompanied by a versioned
\code{secrets.example.dart} template. For a production deployment, we recommend
\textbf{moving the sensitive call behind a backend proxy} so that no secret is
ever embedded in the distributed binary, and injecting configuration via
\code{--dart-define} at build time.

\subsubsection{Principle of least privilege}
The application requests only the \code{INTERNET} and location permissions,
stores no personally identifiable information, and keeps the tile database in the
application's private directory, removing the need for external-storage
permissions.

\begin{figure}[ht]
\centering
\fbox{\begin{minipage}{0.86\textwidth}\centering\vspace{6pt}
\textbf{[ ARCHITECTURE DIAGRAM --- PLACEHOLDER ]}\\[10pt]
\fbox{Presentation Layer: Screens \& Widgets (Home, Map, Assistant\dots)}\\[4pt]
$\downarrow$\\[4pt]
\fbox{State Layer: \code{setState} + \code{IndexedStack}}\\[4pt]
$\downarrow$\\[4pt]
\fbox{Service Layer: ATM \textbullet{} Chat \textbullet{} Location \textbullet{} OfflineMap}\\[4pt]
$\downarrow$\\[4pt]
\fbox{Data Sources: Overpass API \textbullet{} Groq API \textbullet{} MBTiles/SQLite \textbullet{} GPS}
\vspace{6pt}
\end{minipage}}
\caption{Layered architecture of \textit{ATB ATM Locator} (schematic).}
\label{fig:arch-en}
\end{figure}

\section{Implementation Challenges}
\subsection{GPS-induced battery drain}
Continuous position acquisition rapidly depletes the battery. The chosen strategy
favours a \textbf{one-shot acquisition} (\code{getCurrentPosition}) over a
permanent stream, since the static nature of ATM data does not warrant real-time
tracking. This trade-off reduces the energy footprint while preserving high
accuracy at the moment distances are computed.

\subsection{Marker precision and legibility}
Overlaying many markers over a dense area raises a legibility problem. The
current design visually distinguishes selection (enlargement and amber outline)
and keeps the marker count controlled through composable filtering. Marker
\textbf{clustering} is the natural evolution for higher densities.

\subsection{Compliance with the OSM usage policy}
Bulk tile downloading must respect the OSM fair-use policy. The implementation
applies a \textbf{50\,ms throttle} between tiles and a \textbf{User-Agent} that
explicitly identifies the application, preventing any server-side blocking.

\subsection{Resumable downloads}
A network drop must not invalidate a partial download. Each tile already present
in the database is detected and skipped: the download resumes exactly where it
was interrupted, thanks to tile-level \emph{checkpointing}.

\subsection{Coordinate conversion and asynchronous lifecycle}
The conversion between the XYZ ordinate and the \textbf{TMS} ordinate ($y$-axis
inversion) is centralized to avoid alignment errors. Furthermore, asynchronous
callbacks systematically check the widget's \code{mounted} state, and a
single-start flag prevents a download from being re-triggered on every
\code{StatefulBuilder} rebuild.

\section{Performance and CI/CD}
\subsection{Rendering optimization and memory management}
Several levers are used: \code{IndexedStack} avoids tab rebuilding; \code{const}
constructors reduce superfluous rebuilds; the tile database is opened
\textbf{read-only} for browsing and explicitly \textbf{closed} in \code{dispose};
controllers (map, scroll, animation) are released to prevent memory leaks.
Storing tiles as BLOBs in SQLite avoids file-system fragmentation.

\subsection{Continuous integration and deployment pipeline}
The recommended delivery pipeline is built around \textbf{GitHub Actions}:
\begin{itemize}[leftmargin=1.4em]
  \item \textbf{Static analysis}: \code{flutter analyze} with \code{flutter\_lints}
  on every pull request.
  \item \textbf{Testing}: running \code{flutter test} (widgets and service
  logic).
  \item \textbf{Building}: \code{flutter build apk\,/\,appbundle} and \code{ipa},
  with secrets injected via \code{--dart-define}.
  \item \textbf{Signing and publishing}: automation through \textbf{fastlane} to
  test tracks, followed by a \emph{staged rollout}.
\end{itemize}

\section{Conclusion}
The architecture of \textbf{ATB ATM Locator} demonstrates that a clean layered
separation, combined with systematic graceful degradation, is sufficient to
reach production-grade robustness without resorting to costly proprietary
services. \textbf{Sustainability} rests on the exclusive use of open building
blocks (OSM, Overpass), and \textbf{longevity} on the modularity of the service
layer, which will allow substituting a transactional backend, introducing
\code{Riverpod} for global state, and adding biometric authentication and
\emph{push} notifications. The evolution path is thus clearly marked out, which
qualifies the architecture as \emph{future-proof} at the scale of a
consumer-facing banking product.

\end{document}
